\section*{Bab \ref{ch:b2:palaeoenvironments} --- Fosil dan Paleolingkungan}

\begin{solution}{exo:b2:palaeoenvironments:1}
(a) $112/135 = 0.830$ dan $112/150 = 0.747$. (b) Cangkang mati di dasar berkumpul selama bertahun-tahun
atau lebih: cangkang itu mencakup spesies yang dulu hidup di sana, atau jarang hidup di sana, atau terbawa masuk, sehingga
\omterm{def:b2:palaeoenvironments:assemblages}{kumpulan mati} merata-ratakan banyak generasi komunitas. (c) Kelompok bertubuh lunak, khususnya
pemakan detritus dan karnivor seperti cacing: tanpa kerangka, organisme itu membusuk sebelum tertimbun.
\end{solution}

\begin{solution}{exo:b2:palaeoenvironments:2}
(a) $5 \times 165 = \qty{825}{ka}$, sekitar \qty{0.8}{Ma}, di dalam rentang 0.7 sampai 1 juta tahun.
(b) $183.73 - 182.77 = \qty{0.96}{Ma}$, dan $0.96/0.165 = 5.8$: sekitar enam horizon. (c) Suatu
garis keturunan berubah dalam waktu yang jauh kurang dari sejuta tahun; hanya horizon yang memisahkan bentuk-bentuk
yang berurutan, sedangkan zona menggabungkannya.
\end{solution}

\begin{solution}{exo:b2:palaeoenvironments:3}
(a) \omterm{def:b2:palaeoenvironments:ichnofacies}{Iknofasies} Nereites: dasar laut dalam suatu sistem turbidit, di luar tepi paparan, tempat
makanan langka. (b) \omterm{def:b2:palaeoenvironments:ichnofacies}{Iknofasies} Skolithos: pantai depan atau muka pantai. (c) \omterm{def:b2:palaeoenvironments:ichnofacies}{Iknofasies} Cruziana:
lepas pantai, antara dasar gelombang cuaca cerah dan dasar gelombang saat badai. (d) \omterm{def:b2:palaeoenvironments:ichnofacies}{Iknofasies} Zoophycos: pada
Mesozoikum, paparan bawah atau lereng, di bawah dasar gelombang saat badai.
\end{solution}

\begin{solution}{exo:b2:palaeoenvironments:4}
(a) Lapisan A memuat kelompok stenohalin: air laut normal. Lapisan B, yang miskin spesies tetapi kaya
individu, dengan serpihan tumbuhan dan tanpa satu pun kelompok stenohalin, menunjukkan air payau, misalnya
sebuah estuari. (b) Kelompok yang hilang itu banyak, umum, dan terawetkan dengan baik dalam lapisan laut
yang seumur: ketiadaan bersama mereka adalah suatu pola, sedangkan satu fosil dapat hilang karena kebetulan atau
karena pengawetan yang buruk. (c) Sedimennya (saluran pasang surut atau sungai, selaput lumpur, akar), \omterm{def:b2:palaeoenvironments:fossil-kinds}{fosil
jejaknya} (\omterm{def:b2:palaeoenvironments:ichnofacies}{iknofasies} Cruziana dengan sedikit spesies), dan peralihan lapisan B, secara lateral, menjadi
endapan laut atau endapan sungai.
\end{solution}

\begin{solution}{exo:b2:palaeoenvironments:5}
(a) $P = 1/2$ bila $kh/s = \ln 2$: $s = 0.05 \times 0.1/0.693 = \qty{0.0072}{m.yr^{-1}}$,
\qty{7.2}{mm} per tahun. (b) Membagi dua $s$ menggandakan eksponennya: $P = (1/2)^2 = 0.25$, separuhnya.
(c) $k$ dibagi dua: $P = e^{-2.5} = 0.082$ alih-alih $e^{-5} = 0.0067$, $e^{2.5} = 12$ kali lebih
mungkin. (d) Cangkang yang tebal dan kokoh terwakili berlebihan, dan makin berlebihan makin lambat penimbunannya.
\end{solution}

\begin{solution}{exo:b2:palaeoenvironments:6}
(a) Batas atasnya $183.73 + 0.35 = \qty{184.08}{Ma}$: \qty{184.2}{Ma} pada bagan berada
di luarnya, sebesar \qty{0.12}{Ma}. (b) $183.73 - 182.77 = \qty{0.96}{Ma}$. (c) $6/0.96 =
\qty{6.25}{m}$ per juta tahun. (d) Biozona ditentukan oleh fosil, yang kemunculan dan
kepunahannya hampir serentak di wilayah yang luas; ketebalan batuan yang diendapkan selama
waktu itu bergantung pada laju akumulasi setempat.
\end{solution}

\begin{solution}{exo:b2:palaeoenvironments:7}
(a) $0.49\,e^{-1046/3704} = 0.369$ dan $0.49\,e^{-1166/3704} = 0.358$; rata-rata $0.364$. (b) $120
\times (1 - 0.364) = \qty{76.3}{m}$. (c) $76.3/(1 - 0.49) = \qty{149.6}{m}$. (d) Faktor itu
mengandaikan seluruh lapisan yang dipulihkan memiliki porositas permukaan $\varphi_0$; tetapi lapisan
setebal \qty{150}{m} di permukaan kurang berpori di alasnya daripada di puncaknya, sehingga butiran yang sama
mengisi kurang dari \qty{149.6}{m}.
\end{solution}

\begin{solution}{exo:b2:palaeoenvironments:8}
(a) Satu satuan volume memuat $\varphi$ air dan $1 - \varphi$ butiran: $\rho = \varphi
\rho_{\text{air}} + (1 - \varphi)\rho_{\text{butir}}$, yang memberikan rumusnya; $(2710 -
2573)/(2710 - 1000) = 0.080$. (b) $0.51\,e^{-956/4545} = 0.413$. (c) Sementasi: kalsit yang
mengendap di dalam pori, yang dialami karbonat selain kompaksi, menghilangkan porositas jauh lebih banyak
daripada penimbunan saja.
\end{solution}

\begin{solution}{exo:b2:palaeoenvironments:9}
(a) Turbidit laut dalam (\omterm{def:b2:palaeoenvironments:ichnofacies}{iknofasies} Nereites); dasar Jura di paparan bawah atau lereng
(Zoophycos); dasar lepas pantai di antara kedua dasar gelombang (Cruziana); muka pantai (Skolithos);
rawa di darat. (b) Air makin dangkal ke atas: pantai bergerak ke arah laut. (c) Turunnya muka laut,
atau pasokan sedimen yang lebih cepat daripada penurunan dasar, seperti di depan delta yang tumbuh; penurunan muka laut akan
ditemukan pada waktu yang sama di cekungan lain, sedangkan delta bersifat lokal dan sedimennya datang
dari satu sisi. (d) Dalam suksesi yang selaras, lingkungan yang bertetangga saling menyusul;
lompatan melewati lingkungan perantara menunjukkan adanya rumpang waktu atau sesar.
\end{solution}

\begin{solution}{exo:b2:palaeoenvironments:10}
(a) Waktu sampai tertimbun adalah $h/s$ = 100, 10, dan \qty{1}{yr}. Tubuh lunak: $e^{-100} \approx
\num{4e-44}$, $e^{-10} = \num{4.5e-5}$, $e^{-1} = 0.37$. Cangkang ($k = \qty{0.01}{yr^{-1}}$):
$e^{-1} = 0.37$, $e^{-0.1} = 0.90$, $e^{-0.01} = 0.99$. (b) $s = kh/\ln 100 = 0.1/4.605 =
\qty{0.022}{m.yr^{-1}}$, \qty{22}{mm} per tahun. (c) Hampir tidak mungkin, selama berbulan-bulan; tetapi lapisan badai,
turbidit, atau jatuhan abu menimbun dasar dalam hitungan jam. (d) Laju sedimentasi yang tetap: waktu
sampai tertimbun lalu jauh lebih singkat daripada $h/s$ dengan laju rata-rata, dan $P = e^{-kT}$, dengan $T$
waktu sebenarnya sampai tertimbun, mendekati satu bila $T$ jauh lebih kecil daripada $1/k$.
\end{solution}

\begin{solution}{exo:b2:palaeoenvironments:11}
(a) Lapisan tipis $\mathrm dz$ pada kedalaman $z$ memuat butiran $(1 - \varphi(z))\,\mathrm dz$;
di permukaan butiran itu mengisi $\mathrm dz_0$ dengan $(1 - \varphi_0)\,\mathrm dz_0 = (1 -
\varphi(z))\,\mathrm dz$. (b) $\varphi = 0.63\,e^{-816/1960} = 0.415$; $(1 - 0.415)/0.37 =
1.581$, jadi \qty{158.1}{m}. (c) \qty{9.5}{m} terlalu banyak: seperti pada \cref{exo:b2:palaeoenvironments:7},
lapisan yang dipulihkan tidak seluruhnya seberpori permukaannya. (d) Laju butiran,
$h(1-\varphi)/\Delta t$, hanya memakai ketebalan dan porositas masa kini, yang keduanya diukur.
\end{solution}

\begin{solution}{exo:b2:palaeoenvironments:12}
(a) $9.5/0.165 = 57.6$ horizon. (b) $180/57.6 = \qty{3.1}{m}$ batu gamping masa kini, $338/57.6 =
\qty{5.9}{m}$ sedimen segar. (c) Paling banyak \qty{1.5}{m}: $180/1.5 = 120$ selang, 121
sampel. (d) Jeda pengendapan, lapisan yang tidak mengandung amonit, dan horizon yang terkondensasi menjadi
beberapa sentimeter di tempat sedimen datang dengan lambat.
\end{solution}

\begin{solution}{pb:b2:palaeoenvironments:1}
\textbf{1.} $866 \sin \ang{12} = 866 \times 0.2079 = \qty{180.0}{m}$.
\textbf{2.} $481 \times 0.2079 = \qty{100.0}{m}$.
\textbf{3.} Batu pasir, berdasarkan superposisi: sayap-sayapnya hanya miring \ang{12}, sehingga lapisan-lapisannya tidak
terbalik dan yang terbawah adalah yang tertua.
\textbf{4.} $866 - 766 = \qty{100}{m}$ dan $1046 - 866 = \qty{180}{m}$.
\textbf{5.} Sedikitnya \qty{766}{m} di sana; puncak aslinya tererosi, sehingga ketebalan penuhnya
tidak diketahui.
\textbf{6.} 0.70 di permukaan; $0.70\,e^{-766/1408} = 0.406$ pada \qty{766}{m}.
\textbf{7.} Semuanya kelompok stenohalin: airnya bersalinitas laut terbuka.
\textbf{8.} \omterm{def:b2:palaeoenvironments:ichnofacies}{Iknofasies} Cruziana: lepas pantai, antara dasar gelombang cuaca cerah dan dasar gelombang
saat badai.
\textbf{9.} Badai: gelombangnya mencapai dasar, mematahkan cangkang, dan mengendapkannya dalam lapisan
beralas erosif.
\textbf{10.} Air dasarnya miskin oksigen: penggali liang sedikit, lamina tetap terjaga, dan pirit terbentuk.
\textbf{11.} \omterm{def:b2:palaeoenvironments:ichnofacies}{Iknofasies} Zoophycos, pada Mesozoikum di paparan bawah atau lereng, di bawah
dasar gelombang saat badai; kadang-kadang cukup oksigen bagi penggali liang.
\textbf{12.} Tidak: amonit berenang di air di atasnya; cangkangnya tenggelam atau hanyut ke dasar
setelah mati dan membentuk \omterm{def:b2:palaeoenvironments:assemblages}{kumpulan mati}.
\textbf{13.} $184.2 - 174.7 = \qty{9.5}{Ma}$.
\textbf{14.} $180/9.5 = \qty{18.9}{m}$ per juta tahun.
\textbf{15.} $(2710 - 2573)/(2710 - 1000) = 0.080$.
\textbf{16.} $180 \times (1 - 0.080)/(1 - 0.51) = \qty{338}{m}$.
\textbf{17.} $338/9.5 = \qty{35.6}{m}$ per juta tahun.
\textbf{18.} $180 \times 0.92/9.5 = \qty{17.4}{m}$ per juta tahun.
\textbf{19.} Serpih memberikan 11.3 dan \qty{4.4}{m} per juta tahun: laju butiran turun
sekitar empat kali. Karbonat dari gosong dan paparan berhenti datang ketika dasar
makin dalam, dan hanya hujan lumpur yang lambat yang mencapainya.
\textbf{20.} Laut maju: sungai, lalu gosong, dasar lepas pantai, dasar di bawah dasar gelombang saat
badai, lalu laut terbuka tempat kapur terbentuk. Kolom itu merekam suatu transgresi.
\textbf{21.} $57 - 41 = \qty{16}{m}$ selama pengendapan batu gamping, $86 - 57 = \qty{29}{m}$ selama pengendapan serpih,
\qty{45}{m} seluruhnya.
\textbf{22.} $(60 - 5) + 338 + 148.6 = \qty{542}{m}$, dalam $184.2 - 161.5 = \qty{22.7}{Ma}$:
\qty{23.9}{m} per juta tahun.
\textbf{23.} $45/542 = 0.083$: sekitar \qty{8}{\percent}, dan lebih kecil lagi jika diagram siklus membesar-besarkan
amplitudonya. Sisanya berasal dari penurunan dasar.
\textbf{24.} Lumpur gamping batu gamping Varden terakumulasi dengan laju sekitar
\textbf{\qty[text-series-to-math=true]{36}{m}} per juta tahun ketika diendapkan.
\end{solution}
